% 分类目录：dl
\documentclass[12pt, a4paper]{article}
\usepackage[margin=2.5cm]{geometry}
\usepackage{ctex}
\usepackage{amsmath, amssymb}
\usepackage{listings}
\usepackage{xcolor}
\usepackage{tabularx}

\lstset{
    language=Python,
    basicstyle=\ttfamily\small,
    keywordstyle=\color{blue},
    commentstyle=\color{green!50!black},
    stringstyle=\color{red},
    showstringspaces=false,
    breaklines=true,
    numbers=left,
    numberstyle=\tiny\color{gray},
    frame=single
}

\title{知识点：信息论基础 (Information Theory)}
\author{}
\date{}

\begin{document}
\maketitle

\section{核心定义}
\textbf{中文定义}：信息论是一门利用数学统计方法研究信息的量化、存储和通信的学科。在深度学习中，它提供了衡量概率分布差异（如 KL 散度）、评估模型不确定性（如信息熵）及优化目标函数（如交叉熵）的数学骨架。

\textbf{英文定义}：Information Theory is the mathematical study of the quantification, storage, and communication of information. In deep learning, it provides the framework for measuring differences between probability distributions (e.g., KL divergence) and optimizing objective functions.

\section{核心衡量指标}

\subsection{自信息 (Self-Information)}
衡量一个事件发生的惊讶程度。事件 $x$ 的自信息为：
$$ I(x) = -\log p(x) $$
概率越小的事件，包含的信息量越大。

\subsection{信息熵 (Entropy)}
所有可能事件自信息的期望值，衡量随机变量的不确定性：
$$ H(X) = \mathbb{E}_{x \sim P} [I(x)] = -\sum_{x \in \mathcal{X}} p(x) \log p(x) $$

\subsection{相对熵 / KL 散度 (Relative Entropy / KL Divergence)}
衡量两个分布 $P$ 和 $Q$ 之间的差异：
$$ D_{KL}(P || Q) = \mathbb{E}_{x \sim P} \left[ \log \frac{p(x)}{q(x)} \right] = \sum p(x) \log \frac{p(x)}{q(x)} $$
\textbf{注意}：KL 散度是非对称的（即 $D_{KL}(P||Q) \neq D_{KL}(Q||P)$），且始终 $\geq 0$。

\subsection{交叉熵 (Cross-Entropy)}
在深度学习分类任务中最为常用：
$$ H(P, Q) = -\mathbb{E}_{x \sim P} [\log q(x)] = -\sum p(x) \log q(x) $$
它可以拆解为：$H(P, Q) = H(P) + D_{KL}(P || Q)$。
当 $P$（真实标签）固定时，最小化交叉熵等价于最小化 $P, Q$ 之间的 KL 散度。

\subsection{互信息 (Mutual Information)}
衡量两个变量之间的相互依赖程度：
$$ I(X; Y) = H(X) - H(X|Y) = H(Y) - H(Y|X) $$

\section{信息论与深度学习的关系}
\begin{itemize}
    \item \textbf{最大似然估计 (MLE)}：在监督学习中，最小化交叉熵等价于最大化似然函数。
    \item \textbf{VAE 优化}：VAE 的目标函数中包含 KL 散度项，用于约束隐空间分布。
    \item \textbf{信息瓶颈 (Information Bottleneck)}：一种解释深度网络学习过程的理论，认为网络在压缩输入信息的同时保留与输出最相关的特征。
\end{itemize}

\section{概念逻辑关系图}
\begin{verbatim}
[ 联合熵 H(X,Y) ]
      /      \
[ H(X) ] -- [ I(X;Y) ] -- [ H(Y) ]
      \      /
     [ H(X|Y) ]
\end{verbatim}

\section{代码片段 (熵与 KL 散度计算)}
\begin{lstlisting}
import torch
import torch.nn.functional as F

# 1. 计算交叉熵
logits = torch.tensor([[1.0, 2.0, 0.5]])
target = torch.tensor([1]) # 真实类别的索引
loss_ce = F.cross_entropy(logits, target)

# 2. 手动计算 KL 散度
p = torch.tensor([0.1, 0.9])
q = torch.tensor([0.2, 0.8])
kl_div = torch.sum(p * (torch.log(p) - torch.log(q)))

# 3. 使用 PyTorch 内置函数 (注意输入需要 log 空间)
kl_loss = F.kl_div(q.log(), p, reduction='batchmean')
\end{lstlisting}

\section{深度面试题与解答}
\subsection{基础概念}
\textbf{问：为什么 KL 散度不能被称为“距离”？}\\
答：因为“距离”必须满足\textbf{对称性}和\textbf{三角不等式}。KL 散度不满足对称性（$P$ 到 $Q$ 的差异不等于 $Q$ 到 $P$），也不满足三角不等式。它更准确的称呼是“信息增益”或“相对熵”。

\subsection{进阶原理}
\textbf{问：为什么在 GAN 中使用 JS 散度而不是 KL 散度？}\\
答：原生的 GAN 论文发现当 $P$ 和 $Q$ 几乎不重叠时，KL 散度是常数，导致梯度反传失败。JS 散度是 KL 散度的对称变体，取值在 $[0, \log 2]$ 之间，在分布不重叠时依然能提供更好的数值特性。

\subsection{工程实战}
\textbf{问：Label Smoothing（标签平滑）是如何从信息论角度解释的？}\\
答：标签平滑是将 One-hot 的 $P$ 变成一个更加熵增的分布。从交叉熵角度看，它限制了模型对样本的“极端自信”，增加了输出分布的熵值，从而起到正则化作用，防止过拟合。

\section{指标对比总结表}

\begin{table}[h]
\centering
\begin{tabularx}{\textwidth}{|l|X|l|}
\hline
\textbf{指标} & \textbf{核心用途} & \textbf{取值范围} \\
\hline
信息熵 & 衡量单一变量的不确定性 & $[0, \infty)$ \\
\hline
交叉熵 & 监督学习的分类损失函数 & $[0, \infty)$ \\
\hline
KL 散度 & 衡量分布间差异（生成模型常用） & $[0, \infty)$ \\
\hline
互信息 & 特征选择、表征学习相关性衡量 & $[0, \min(H(X),H(Y))]$ \\
\hline
\end{tabularx}
\end{table}

\section{参考文献与拓展}
\begin{enumerate}
    \item \textbf{奠基之作}: Shannon, C. E. (1948). \textit{A Mathematical Theory of Communication}.
    \item \textbf{经典教科书}: Cover, T. M., & Thomas, J. A. \textit{Elements of Information Theory}.
    \item \textbf{DL 应用}: Goodfellow, I., et al. \textit{Deep Learning} (Chapter 3: Probability and Information Theory).
\end{enumerate}

\end{document}
